\chapter{Path Integrals}

Path integrals were invented as an alternative but equivalent quantization procedure to canonical quantization. In this chapter, we will introduce the path integral formulation of quantum field theory and show how it can be used to compute correlation functions and scattering amplitudes.

Applying path integrals to traditional quantum mechanics seems a bit complicated and unnecessary, but it becomes very useful in quantum field theory, where the canonical quantization procedure can be quite cumbersome. Moreover, path integrals provide a more intuitive picture of quantum mechanics and quantum field theory, as they allow us to think of particles as taking all possible paths between two points, rather than just a single path.

Among other things, path integrals are particularly useful:
\begin{itemize}
    \item \textbf{Lorentz invariance is manifest}: the path integral formulation is manifestly Lorentz invariant, since it is a lagrangian formulation of the theory, and thus it does not require a choice of time coordinate or a specific reference frame (we do not need the hamiltonian).
    \item \textbf{Quantization of gauge field theories} is more straightforward: the path integral formulation allows us to easily implement gauge fixing and Faddeev-Popov ghosts, which are necessary for quantizing gauge theories.
    \item \textbf{Numerical simulations for strongly interacting theories} (where perturbation theory fails): path integrals can be evaluated numerically using Monte Carlo methods, which is particularly useful for studying strongly interacting theories such as quantum chromodynamics (QCD) on the lattice (the idea is to discretize spacetime into a lattice and then evaluate the path integral numerically).
    \item The use of \textbf{functional integrals} allows us to easily derive the generating functional for correlation functions, which is a powerful tool for computing correlation functions and scattering amplitudes in quantum field theory. For examples \textit{symmetries} in QFT are much easier to understand in the path integral formulation, since they can be implemented as transformations of the fields in the path integral, and thus we can easily derive Ward identities and other relations between correlation functions.
\end{itemize}

\section{Path Integral Formulation of Quantum Mechanics}

The idea of the path integral formulation is to express the transition amplitude between two states as a sum over all possible paths that the system can take between those states. In quantum mechanics, the most used eexample to illustrate the path integral formulation is the \textbf{double slit experiment}, where we can think of the particle as taking all possible paths from the source to the detector, and the interference pattern arises from the constructive and destructive interference of these paths.

    [\textit{drawing of source of particles, slits and detector}]\TODO{Add drawing of double slit experiment}

\[
    A(x \to y) = A (x \to a \to y) + A (x \to b \to y),
\]
where \(a\) and \(b\) are the two slits. In the path integral formulation, we can express the transition amplitude as a sum over all possible paths that the particle can take from \(x\) to \(y\), but we can do more, adding virtually an infinite number of screens with slits in between \(x\) and \(y\), and thus we can express the transition amplitude as a sum over all possible paths that the particle can take from \(x\) to \(y\) through any number of slits. In the limit of infinitely many screens with infinitely many slits, we can express the transition amplitude as a sum over all possible paths that the particle can take from \(x\) to \(y\), without any constraints.

    [\textit{draw of infinite walls and slits among x and y}]\TODO{Add drawing of infinite walls and slits among x and y}

If \(\gamma\) is a specific path from \(x\) to \(y\), then we can associate to it a complex number \(A(\gamma)\) called the amplitude of the path, and then we can express the transition amplitude as a sum over all possible paths:
\[
    A(x \to y) = \sum_{\gamma} A(x \to y,\, \gamma).
\]

In order to implement this idea, we can discretize the time interval between \(x\) and \(y\) into \(N\) small intervals of size \(\epsilon\), and then we can express the transition amplitude as a sum over all possible paths that the particle can take from \(x\) to \(y\) through these small intervals:
\[
    \bra{f;\,t_f}\ket{i;\,t_i}, \quad t_i < t_1 < \cdots < t_n < t_f, \quad t_{k+1} - t_k = \epsilon,
\]
where we have inserted \(N-1\) complete sets of states at intermediate times \(t_1, \ldots, t_{N-1}\). Then we can express the transition amplitude as a sum over all possible paths that the particle can take from \(x\) to \(y\) through these small intervals by adding the resolution of the identity at each intermediate time:
\[
    \bra{f;\,t_f}\ket{i;\,t_i} = \int \mathrm{d} q_1 \cdots \mathrm{d}q_n \bra{f;\,t_f}\ket{q_{n};\,t_{n}} \bra{q_{n};\,t_{n}}\ket{q_{N-1};\,t_{N-1}} \cdots \bra{q_1;\,t_1}\ket{i;\,t_i}.
\]
The result to this expression is
\[
    \bra{f;\,t_f}\ket{i;\,t_i} = N \int_{q(t_i) = x_i}^{q(t_f) = x_f} \mathcal{D} q(t) e^{i S[q(t)]},
\]
where \(S[q(t)] = \int \mathrm{d}t \, L(q(t), \dot{q}(t))\) is the action of the path \(q(t)\), and \(\mathcal{D} q(t)\) is a measure that assigns a weight to each path. The normalization factor \(N\) is chosen such that the transition amplitude is properly normalized. In order to understand this expression more intuitively, we can always go back to its discretized version:
\[
    \bra{f;\,t_f}\ket{i;\,t_i} = N \int \mathrm{d} q_1 \cdots \mathrm{d}q_n \exp[i \sum_{k} L(q_k, \dot{q}_{k}) \cdot \epsilon], \quad \dot{q}_k = \frac{q_{k+1} - q_k}{\epsilon}.
\]

We can also generalize the last expression using the time ordering operator \(\hat{T}\):
\[
    \bra{f;\,t_f} \hat{T} \{\hat{q}(t_a) \hat{q}(t_b) \cdots\} \ket{i;\,t_i} = N \int_{q(t_i) = x_i}^{q(t_f) = x_f} \mathcal{D} q(t) e^{i S[q(t)]} q(t_a) q(t_b) \cdots, \quad t_i < t_a < t_b < \cdots < t_f,
\]
where the time ordering operator \(\hat{T}\) ensures that the operators are ordered in time from right to left, with the earliest time on the right and the latest time on the left. This expression allows us to compute correlation functions of the form \(\bra{f;\,t_f} \hat{T} \{\hat{q}(t_a) \hat{q}(t_b) \cdots\} \ket{i;\,t_i}\) using path integrals.

There exists also a discretized version of the last expression, which is more intuitive and easier to understand:
\[
    \bra{f;\,t_f} \hat{T} \{\hat{q}(t_a) \hat{q}(t_b) \cdots\} \ket{i;\,t_i} = N \int \mathrm{d} q_1 \cdots \mathrm{d}q_n \exp[i \sum_{k} L(q_k, \dot{q}_{k}) \cdot \epsilon] q(t_a) q(t_b) \cdots,
\]
where we have replaced the operators \(\hat{q}(t_a)\) and \(\hat{q}(t_b)\) with the corresponding functions \(q(t_a)\) and \(q(t_b)\) in the path integral as did before. This expression allows us to compute correlation functions of the form \(\bra{f;\,t_f} \hat{T} \{\hat{q}(t_a) \hat{q}(t_b) \cdots\} \ket{i;\,t_i}\) using path integrals in a more intuitive way. Time ordering is basically implemented for free, since the path integral already sums over all possible paths, and thus it automatically takes into account all possible time orderings of the operators. \QUESTION{time ordering.}

\subsection{Classical Limit}

In the classical limit, the path integral is dominated by the path that minimizes the action, which is the classical path. This is because in the classical limit, the action is much larger than \(\hbar \to 0\), and thus the contributions from paths that deviate from the classical path are suppressed by a factor of \(e^{i S/\hbar}\),\footnote{We omitted the \(\hbar\) factor for simplicity, but via dimensional analysis it can be easily recovered.} which oscillates rapidly and averages out to zero. Therefore, in the classical limit, we can approximate the path integral by evaluating it at the classical path, which gives us the classical equations of motion.

Thus the most important contribution to the path integral comes from the classical path, which is the path that minimizes the action. This is known as the principle of least action, and it is a fundamental principle in classical mechanics that states that the motion of a system is determined by the path that minimizes the action. From these procedure we can easily compute the Euler-Lagrange equations of motion, which are the equations that govern the dynamics of the system in classical mechanics. In the path integral formulation, we can derive the Euler-Lagrange equations by performing a variation of the action with respect to the path, and then setting the variation to zero.

\section{Path Integral Formulation of QFT}

The path integral formulation can be generalized with the use of some handy substitutions:
\[
    \begin{dcases}
        q \to \phi(x), \\
        t, \, \mathrm{d} t \to x^{\mu}, \, \mathrm{d}^4 x,
    \end{dcases}
\]

... green's finctions... \(\ket{i} = \ket{f} = \ket{\Omega}\)...

\[
    \bra{\Omega} \hat{T} \{\hat{\phi}(x_1) \hat{\phi}(x_2) \cdots \hat{\phi} (x_n)\} \ket{\Omega} = N \int \mathcal{D} \phi(x) e^{i S[\phi(x)]} \phi(x_1) \phi(x_2) \cdots \phi(x_n),
\]
where now the action is \(S[\phi(x)] = \int \mathrm{d}^4 x \, \mathcal{L}(\phi(x), \partial_\mu \phi(x))\). The normalization \(N\) can be fixed by requiring that the vacuum-to-vacuum transition amplitude is equal to one \(\bra{\Omega} \ket{\Omega} = 1\):
\[
    N = \frac{1}{\int \mathcal{D} \phi(x) e^{i S[\phi(x)]}}.
\]


In order to get the discretized version of the last expression, we can discretize spacetime into a lattice of points, \(x_i \to \phi_i \equiv \phi(x_i)\) and

    [\textit{draw of lattice}]\TODO{Add drawing of lattice}

\[
    \bra{\Omega} \hat{T} \{\hat{\phi}(x_1) \hat{\phi}(x_2) \cdots \hat{\phi} (x_n)\} \ket{\Omega} = N \int \mathrm{d} \phi_1 \cdots \mathrm{d} \phi_N \exp[i \sum_{k=1}^N \Delta^4 x \mathcal{L}(\phi_k)] \phi_1 \phi_2 \cdots \phi_N,
\]

where again... and in the limit for \(n \to  \infty\) ...

But how do we evaluate the path integral? In general, it is very difficult to evaluate the path integral exactly, but we have some options based on the context in which we are working:
\begin{itemize}
    \item \textbf{Lattice QFT}: we can evaluate the path integral numerically using Monte Carlo methods, which is particularly useful for studying strongly interacting theories such as quantum chromodynamics (QCD) on the lattice (the idea is to discretize spacetime into a lattice and then evaluate the path integral numerically). In the end we could infere the properties at the continuum limit by taking the limit of the lattice spacing going to zero.
    \item \textbf{Perturbation Theory}: if the lagrangian of the theory is quadratic in the fields, then we can evaluate the path integral exactly using \textit{gaussian integrals}. If the lagrangian is not quadratic, then we can treat the non-quadratic terms as perturbations and evaluate the path integral perturbatively using Feynman diagrams: we can Taylor expand the exponential of the action around the free theory, and then we can use Wick's theorem to evaluate the resulting correlation functions in terms of the free propagators. This is the standard approach for computing scattering amplitudes in quantum field theory, and it is based on the idea that we can treat interactions as small perturbations to the free theory.
\end{itemize}

\subsection{Gaussian Integrals}

One dimensional gaussian integrals are given by
\[
    \int_{-\infty}^{\infty} \mathrm{d} z e^{-\tfrac{a}{2} z^2} = \sqrt{\frac{2\pi}{a}}, \quad a > 0.
\]
Now the generalization to \(n\) dimensions is given by \(\mathbf{z} = (z_1, z_2, \ldots, z_n)\)
\[
    \int \mathrm{d} z_1 \cdots \mathrm{d} z_n \; \exp[-\frac{1}{2} \mathbf{z}^T \hat{A} \mathbf{z} + \mathbf{J} \cdot \mathbf{z}] = \sqrt{\frac{(2\pi)^{n/2}}{\sqrt{\det \hat{A}}}} \exp[\frac{1}{2} \mathbf{J}^T \hat{A}^{-1} \mathbf{J}].
\]

As we said cubic or higher order terms in the lagrangian can be treated as perturbations, and thus we can evaluate the path integral perturbatively using Feynman diagrams: we can Taylor expand the exponential of the action around the free theory
\[
    \int_{-\infty}^{\infty} \mathrm{d} z \exp[-\frac{1}{2}] \dots \dots
\]
and
\[
    \dots
\]


\subsection{Computing Green's Functions}

Since we are working in a space of functions, we have to define the functional derivative as
\[
    \frac{\partial }{\partial J(x)} J(y) = \delta^{(4)}(x-y),
\]
which implies
\[
    \frac{\partial }{\partial J(x)} \int \mathrm{d}^4 y \, J(y) \phi(y) = \phi(x),
\]
which is the continuum version of the discrete case \(\frac{\partial }{\partial J_i} \sum_j J_j \phi_j = \phi_i\) (and \(\frac{\partial}{\partial J_i} J_j = \delta_{ij}\)). This functional derivative obeys the same properties as the ordinary derivative, such as linearity and the product rule. We can use this functional derivative to compute correlation functions by introducing a source term \(J(x)\) in the action, and then taking functional derivatives with respect to \(J(x)\) to obtain correlation functions of the fields.

Since we can almost treat it as a normal derivative we have
\[
    \frac{\partial}{\partial J(x)} \int \mathrm{d}^4 y \, (\partial_\mu J(y)) \phi(y) = - \partial_\mu \phi(x),
\]
having integrated by parts and assuming that the fields vanish at infinity. This is a very useful property, since it allows us to compute correlation functions involving derivatives of the fields by simply taking functional derivatives with respect to the source term.

\begin{example}[Exercise]
    Consider a lagrangian \(\mathcal{L} = \mathcal{L} (l_k, \partial_\mu l_k)\) show that the equations of motion can ...
\end{example}

We can also define the \textbf{generating functional} \(Z[J]\) as
\[
    Z[J] \equiv \int \mathcal{D} \phi e^{i S[\phi(x)] + i \int \mathrm{d}^4 x J(x) \phi(x)},
\]
where \(J(x)\) can be thought as a classical source that couples to the quantum field \(\phi(x)\). The generating functional \(Z[J]\) is a powerful tool for computing correlation functions, since we can obtain correlation functions by taking functional derivatives of \(Z[J]\) with respect to \(J(x)\) and then setting \(J(x) = 0\):
\[
    \bra{\Omega} \hat{T} \{\hat{\phi}(x_1) \hat{\phi}(x_2) \cdots \hat{\phi} (x_n)\} \ket{\Omega} = \frac{1}{Z[0]} \left. \left(-i\frac{\partial}{\partial J(x_1)}\right) \left(-i\frac{\partial}{\partial J(x_2)}\right) \cdots \left(-i\frac{\partial}{\partial J(x_n)}\right) Z[J] \right|_{J=0},
\]
where the normalization factor \(1/Z[0]\) is necessary to ensure that the correlation functions are properly normalized. The generating functional \(Z[J]\) can be evaluated exactly for free theories, and it can be evaluated perturbatively for interacting theories using Feynman diagrams.

If we can compute the generating functional \(Z[J]\) for a given theory, then we can compute all correlation functions of the theory by taking functional derivatives with respect to \(J(x)\) and then setting \(J(x) = 0\): we basically know everything about the theory, and thus we can compute all scattering amplitudes and other physical observables. This is why the generating functional \(Z[J]\) is such a powerful tool in quantum field theory.

In order to compute exactely the generating functional \(Z[J]\) for a free theory, we can use the fact that the action is quadratic in the fields, and thus we can evaluate the path integral exactly using gaussian integrals: rewriting the action as
\[
    i \int \mathrm{d}^4 x \, \left( \mathcal{L} + \phi J \right) = -\frac{i}{2} \int \mathrm{d}^4 x \phi \left( \partial^2 + m^2 - i \epsilon\right) \phi + i \int \mathrm{d}^4 x \, J \phi,
\]
where we have added a small imaginary part \(i \epsilon\) to the mass term in order to ensure that the path integral converges. Then we can evaluate the path integral using gaussian integrals, and we find
\[
    Z[J] = Z[0] \exp\left[i \int \mathrm{d}^4 x \left(-\frac{1}{2} \phi (\partial^2 + m^2 - i \epsilon) \phi + J \phi \right)\right],
\]
where we can rename \((\partial^2 + m^2 - i \epsilon) = \hat{A}\) as the matrix used when defining the gaussian integral. We have to note a couple of things:
\[
    \begin{dcases}
        \hat{A} = \partial^2 + m^2 - i \epsilon,                                                     \\
        \hat{A} D_F(x-y) = \delta^{(4)}(x-y),                                                        \\
        \sum_j A_{ij} (D_F)_{jk} = \delta_{ik} \to \int \mathrm{d}^4 p \, e^{ip \cdot (x-y)} \Pi(p), \\
        i (\partial^2 + m^2 - i \epsilon) \Pi(p) = 1 \implies \Pi(p) = \frac{i}{p^2 - m^2 + i \epsilon}.
    \end{dcases}
\]

Thus now we can rewrite the generating functional \(Z[J]\) as
\[
    Z[J] = Z_0 \exp\left[-\frac{1}{2} \int \mathrm{d}^4 x \int  \mathrm{d}^4 y J(x) D_F(x-y) J(y)\right],
\]
which is the exact expression for the generating functional of a free theory. From this expression, we can compute all correlation functions of the free theory by taking functional derivatives with respect to \(J(x)\) and then setting \(J(x) = 0\).

\subsubsection{Green's Functions of a Free Theory}

\begin{example}[Two point green's function]
    The two point green's function is given by
    \[
        \bra{0} \hat{T} \{\hat{\phi}(x_1) \hat{\phi}(x_2)\} \ket{0} = \frac{1}{Z[0]} \left. \left(-i\frac{\partial}{\partial J(x_1)}\right) \left(-i\frac{\partial}{\partial J(x_2)}\right) Z[J] \right|_{J=0}.
    \]
    Then we can compute the functional derivatives and we find
    \[
        \begin{aligned}
            \bra{0} \hat{T} \{\hat{\phi}(x_1) \hat{\phi}(x_2)\} \ket{0} & = \frac{1}{Z[0]} \left. \left(-i\frac{\partial}{\partial J(x_1)}\right) \left[-\frac{1}{2} \int \mathrm{d}^4 y D_F(x_2-y) J(y) - \frac{1}{2} \int \mathrm{d}^4 x J(x) D_F(x-x_2) \right] \frac{Z[J]}{Z_0}\right|_{J=0} \\
                                                                        & = \frac{1}{2} D_F(x_2 - x_1) + \frac{1}{2} D_F(x_1 - x_2) = D_F(x_2 - x_1),
        \end{aligned}
    \]
    which is the Feynman propagator of the free theory.
\end{example}

\begin{example}[Four point green's function]
    If we rename for convenience
    \[
        \begin{dcases}
            D_F(x-y) = D_{xy},                                                                                       \\
            J(x) = J_x,                                                                                              \\
            \int \mathrm{d}^4 x \int \mathrm{d}^4 y J(x) D_F(x-y) J(y) = \sum_{x,y} J_x D_{xy} J_y = J_x D_{xy} J_y, \\
            \frac{\partial}{\partial J_z} (J_x D_{xy} J_y) = D_{zy} J_y + J_x D_{xz} = 2 J_x D_{xz},
        \end{dcases}
    \]
    then we can compute the four point green's function as
    \[
        \begin{aligned}
            \bra{0} \hat{T} \{\hat{\phi}_1 \hat{\phi}_2 \hat{\phi}_3 \hat{\phi}_4\} \ket{0} & = \left. (-i)^4 \frac{\partial^4}{\partial J_1 \partial J_2 \partial J_3 \partial J_4} \exp[-\frac{1}{2} J_x D_{xy} J_y] \right|_{J=0} \\
                                                                                            & = \frac{\partial}{\partial J_1} \cdots \frac{\partial}{\partial J_3} (-J_x D_{x4}) \exp[\cdots] \bigg|_{J=0}                           \\
                                                                                            & = \frac{\partial}{\partial J_1} \frac{\partial}{\partial J_2} \left(-D_{34} + J_x D_{x4} J_y D_{yz} \right)\exp[\cdots] \bigg|_{J=0}   \\
                                                                                            & =  \cdots = D_{14} D_{23} + D_{24} D_{13} + D_{12} D_{34}                                                                              \\
                                                                                            & = \sum \text{ full contractions of } \phi_1 \phi_2 \phi_3 \phi_4,
        \end{aligned}
    \]
    which is the sum of all possible contractions of the four fields, and thus it is the sum of all possible Feynman diagrams with four external legs.
\end{example}


\subsection{Perturbation Theory and Feynman Diagrams} \TODO{check}

Until now, we have successfully treated free theories in the path integral formulation by recognizing the free Lagrangian as yielding a functional Gaussian integral, which we were able to evaluate exactly. Let us now check that interacting theories can also be systematically treated within this formulation, and verify that we recover the familiar perturbative expansion.

Consider an interacting theory whose Lagrangian density can be split into a free part and an interaction part: $\mathcal{L} = \mathcal{L}_{\text{free}} + \mathcal{L}_{\text{int}}$.
We want to evaluate the $n$-point Green's function, which is given by the path integral normalized by the partition function of the free theory ($Z_0$):
\[
    \begin{aligned}
        \langle \Omega | T \{ \hat{\phi}(x_1) \cdots \hat{\phi}(x_n) \} | \Omega \rangle \times Z_0 & = \int \mathcal{D}\phi \, \exp\left[ i \int \mathrm{d}^4x (\mathcal{L}_{\text{free}} + \mathcal{L}_{\text{int}}) \right] \phi(x_1) \cdots \phi(x_n)                                         \\
                                                                                                    & = \int \mathcal{D}\phi \, \exp\left[ i \int \mathrm{d}^4x \mathcal{L}_{\text{free}} \right] \exp\left[ i \int \mathrm{d}^4z \mathcal{L}_{\text{int}}(z) \right] \phi(x_1) \cdots \phi(x_n).
    \end{aligned}
\]
Assuming the interaction is weak (controlled by a small coupling constant), we can Taylor-expand the exponential of the interaction term:
\[
    = \int \mathcal{D}\phi \, \exp\left[ i \int \mathrm{d}^4x \mathcal{L}_{\text{free}} \right] \sum_{k=0}^{\infty} \frac{i^k}{k!} \int \mathrm{d}^4z_1 \cdots \mathrm{d}^4z_k \, \mathcal{L}_{\text{int}}(z_1) \cdots \mathcal{L}_{\text{int}}(z_k) \phi(x_1) \cdots \phi(x_n).
\]
Notice that for each term in the sum, the path integral over the exponential of the free action acts as a Gaussian weight for the polynomial function composed of the product of fields $\mathcal{L}_{\text{int}}(z_1) \cdots \mathcal{L}_{\text{int}}(z_k) \phi(x_1) \cdots \phi(x_n)$.
Using the functional equivalent of Wick's theorem (or by taking successive functional derivatives of the free generating functional), this integral evaluates exactly to the sum of all possible full contractions between the fields:
\[
    = \sum_{k=0}^{\infty} \frac{i^k}{k!} \int \mathrm{d}^4z_1 \cdots \mathrm{d}^4z_k \sum \left( \text{all full contractions in } \mathcal{L}_{\text{int}}(z_1) \cdots \mathcal{L}_{\text{int}}(z_k) \phi(x_1) \cdots \phi(x_n) \right).
\]
This expression is exactly what we found when we first computed Green's functions of interacting theories using the canonical operator formalism, where the same conclusion was derived by applying Wick's theorem to the Dyson series $\langle 0 | T \{ \phi(x_1) \cdots \phi(x_n) \exp[i \int \mathrm{d}^4x \mathcal{H}_{\text{int}}] \} | 0 \rangle$.

We already know that the sum over contractions can be efficiently organized and computed using Feynman diagrams and their associated Feynman rules. Moreover, it is straightforward to identify that the normalization factor $Z_0$ in the denominator exactly cancels the disconnected "vacuum bubble" diagrams, meaning that dividing by $Z_0$ amounts to considering only \textbf{connected diagrams}.

We have thus successfully reproduced the Feynman rules starting from the path integral formalism! This proves that the canonical operator formalism and the path integral formalism yield identical results for Green's functions, and therefore (via the LSZ reduction formula) also for scattering amplitudes.

\begin{remark}
    Strictly speaking, we have shown this equivalence explicitly only for interaction terms $\mathcal{L}_{\text{int}}$ that do not contain derivatives. However, as previously discussed regarding the treatment of derivatives in the path integral, this formalism guarantees that we can safely use $i\mathcal{L}_{\text{int}}$ instead of $-i\mathcal{H}_{\text{int}}$ to derive Feynman rules, even when derivative couplings are present, without worrying about generating spurious contact terms from the time-ordering operator.
\end{remark}

\subsubsection{Feynman Rules via Functional Derivatives}

The path integral formalism provides an extremely elegant and direct recipe for deriving Feynman rules for interaction vertices simply by taking functional derivatives of the action $\int \mathrm{d}^4x \, i \mathcal{L}_{\text{int}}$.

In momentum space, we first rewrite each field in the interaction Lagrangian using its Fourier transform, e.g., $\phi(x) = \int \mathrm{d}^4k \, \phi(k) e^{ik \cdot x}$.
Then, the Feynman rule for a specific vertex with incoming and outgoing momenta $p_i$ is obtained by applying the following recipe:
\begin{itemize}
    \item For each \textbf{incoming particle} (or outgoing anti-particle) with momentum $p$, apply the functional derivative $\frac{\delta}{\delta \phi(-p)}$.
    \item For each \textbf{outgoing particle} (or incoming anti-particle) with momentum $p$, apply the functional derivative $\frac{\delta}{\delta \phi^*(+p)}$.
\end{itemize}
After computing all the functional derivatives, you must set all fields to zero ($\phi, \phi^* \to 0$) to isolate the purely numerical coefficient. This procedure will naturally yield the algebraic Feynman rule multiplied by the overall momentum conservation delta function $\delta^{(4)}(p_{\text{in}} - p_{\text{out}})$.

\begin{example}[Interacting Lagrangian in sQED]
    Let us apply this recipe to derive the 3-point vertex rule in scalar QED. The interaction term in the Lagrangian for a complex scalar field coupled to a photon contains a derivative coupling:
    \[
        \mathcal{L}_{\text{int}} \supset -ie A^\mu [ \phi^* (\partial_\mu \phi) - (\partial_\mu \phi^*) \phi ].
    \]
    Consider the vertex where a scalar particle enters with momentum $p_1$, a photon enters with momentum $p_3$ (and Lorentz index $\mu$), and the scalar particle exits with momentum $p_2$.

    First, we Fourier transform the fields inside the action integral $\int \mathrm{d}^4x \, i \mathcal{L}_{\text{int}}$ and convert the coordinate derivatives into momentum factors ($\partial_\mu \to ik_\mu$):
    \[
        \int \mathrm{d}^4x \int \mathrm{d}^4k_1 \mathrm{d}^4k_2 \mathrm{d}^4k_3 \, e^{ix \cdot (k_1 + k_2 + k_3)} \left( e A^\mu(k_3) \phi^*(k_2) \phi(k_1) (ik_{1\mu} - ik_{2\mu}) \right).
    \]
    To find the vertex rule, we apply the appropriate functional derivatives for the incoming scalar (momentum $p_1$), outgoing scalar (momentum $p_2$), and incoming photon (momentum $p_3$, index $\nu$):
    \[
        \frac{\delta}{\delta \phi(-p_1)} \frac{\delta}{\delta \phi^*(p_2)} \frac{\delta}{\delta A^\nu(-p_3)} \left[ \dots \right].
    \]
    We use the fundamental properties of functional derivatives in momentum space:
    \[
        \frac{\delta \phi(k)}{\delta \phi(-p)} = \delta^{(4)}(-p-k), \quad \frac{\delta \phi^*(k)}{\delta \phi^*(p)} = \delta^{(4)}(p-k), \quad \frac{\delta A^\mu(k)}{\delta A^\nu(-p)} = \delta^\mu_\nu \delta^{(4)}(-p-k).
    \]
    Applying these derivatives to our fields $\phi(k_1)$, $\phi^*(k_2)$, and $A^\mu(k_3)$ respectively, we get:
    \[
        \int \mathrm{d}^4x \int \mathrm{d}^4k_1 \mathrm{d}^4k_2 \mathrm{d}^4k_3 \, e^{ix \cdot (k_1 + k_2 + k_3)} \left( e \, \delta^\mu_\nu \delta^{(4)}(-p_3 - k_3) \, \delta^{(4)}(p_2 - k_2) \, \delta^{(4)}(-p_1 - k_1) \, (ik_{1\mu} - ik_{2\mu}) \right).
    \]
    Integrating over the $k_i$ variables simply forces $k_1 = -p_1$, $k_2 = p_2$, and $k_3 = -p_3$:
    \[
        = \int \mathrm{d}^4x \, e^{ix \cdot (-p_1 + p_2 - p_3)} \, e \, \delta^\mu_\nu \left( i(-p_{1\mu}) - i(p_{2\mu}) \right).
    \]
    Finally, integrating over spacetime $x$ produces the overall momentum conservation delta function $(2\pi)^4\delta^{(4)}(-p_1 + p_2 - p_3)$. The algebraic coefficient multiplying this delta function is exactly the Feynman rule for the vertex:
    \[
        \text{Vertex Rule} = e (-i p_{1\nu} - i p_{2\nu}) = -ie (p_1 + p_2)_\nu.
    \]
    \textit{(Exercise: Try applying the same functional derivative procedure to the quadratic 4-point "seagull" interaction term $e^2 A_\mu A^\mu \phi^* \phi$ to derive its Feynman rule $2ie^2 g_{\mu\nu}$).}
\end{example}

\subsection{Fermionic Path Integrals} \TODO{check}

Inside the standard path integral, everything commutes because we are treating quantum operators as ordinary classical functions. To extend this formalism to fermions, which physically anticommute, we must think of fermionic fields inside the path integral as fundamentally "anticommuting" objects. We achieve this formally by introducing \textbf{Grassmann variables} (or Grassmann numbers) \(\theta_i\). In practice, you only need to know that they satisfy the following properties:
\begin{itemize}
    \item \textbf{Anticommutation}: \(\theta_i \theta_j + \theta_j \theta_i = 0\) for all \(i, j\). In particular, setting \(i=j\) implies that the square of any Grassmann variable is exactly zero: \(\theta_i^2 = 0\).
    \item \textbf{Integration and differentiation}: These operations are well-defined and, remarkably, they coincide. The integral of a Grassmann variable acts exactly like its derivative:
          \[
              \int \mathrm{d} \theta_1 \cdots \mathrm{d} \theta_n X = \frac{\partial}{\partial \theta_n} \cdots \frac{\partial}{\partial \theta_1} X.
          \]
          This definition is consistent with the anticommutation property. Since \(\theta_i^2 = 0\), the Taylor expansion of any function of a single Grassmann variable terminates at the linear term (e.g., \(f(\theta) = a + b\theta\)). Consequently, the integral of any product of Grassmann variables is zero unless it contains exactly one factor of each variable. Furthermore, we naturally have
          \[
              \frac{\partial \theta_i}{\partial \theta_j} = \delta_{ij}.
          \]
    \item \textbf{Gaussian integrals}: You can use the above rules to compute a Gaussian integral over multidimensional Grassmann variables:
          \[
              \int \mathrm{d} \theta_1 \mathrm{d} \overline{\theta}_1 \cdots \mathrm{d} \theta_n \mathrm{d} \overline{\theta}_n \exp\left(-\sum_{i,j} \overline{\theta}_i A_{ij} \theta_j + \sum_{j} \left(\overline{\eta}_j \theta_j + \overline{\theta}_j \eta_j\right) \right) = \det A \, \exp\left(\sum_{i,j} \overline{\eta}_i (A^{-1})_{ij} \eta_j\right).
          \]
          Here \(A\) is a matrix, while \(\eta_j\) and \(\overline{\eta}_j\) are additional Grassmann variables that serve as sources for the fermionic fields. This is the fundamental tool for evaluating fermionic path integrals. \textit{Crucial difference:} Note that, unlike the bosonic case, the determinant \(\det A\) appears in the \textbf{numerator} rather than the denominator.
\end{itemize}

We can thus write a path integral for the Dirac field and, more generally, its generating functional:
\[
    Z[\eta, \overline{\eta}] = \int \mathcal{D} \overline{\psi} \mathcal{D} \psi \exp\left[i \int \mathrm{d}^4 x \left( \overline{\psi}(i \slashed{\partial} - m) \psi + \overline{\eta} \psi + \overline{\psi} \eta\right)\right],
\]
where we may shift \(m \to m-i\epsilon\) to ensure formal convergence, just as we did in the bosonic case.

It is important to note that \(\psi\) and \(\overline{\psi}\) must be treated as completely independent variables, and thus each requires its own independent source field, \(\overline{\eta}(x)\) and \(\eta(x)\) respectively. Furthermore, because the Lagrangian scalar \(\mathcal{L}\) must be a commuting object (so that the action \(S\) is a standard commuting number), the source fields \(\eta\) and \(\overline{\eta}\) must themselves be \textbf{anticommuting} Grassmann variables to perfectly cancel the anticommuting nature of \(\psi\) and \(\overline{\psi}\) in the source terms.

Via the same completing-the-square method used for scalar fields (which is left as a highly recommended exercise), we can compute \(Z\) for the free Dirac theory, obtaining:
\[
    Z[\eta, \overline{\eta}] = Z_0 \exp\left[-\int \mathrm{d}^4 x \int \mathrm{d}^4 y \, \overline{\eta}(x) S_F(x-y) \eta(y)\right],
\]
where \(S_F(x-y)\) is the Dirac propagator. It is defined as the inverse of the Dirac operator \(i \slashed{\partial} - m + i \epsilon\):
\[
    S_F(x-y) = \int \frac{\mathrm{d}^4 p}{(2\pi)^4} e^{-i (x-y)\cdot p} \frac{i}{\slashed{p} - m + i \epsilon} = \int \frac{\mathrm{d}^4 p}{(2\pi)^4} e^{-i (x-y)\cdot p} \frac{i (\slashed{p} + m)}{p^2 - m^2 + i \epsilon} = \bra{0} \hat{T} \{\hat{\psi}(x) \overline{\hat{\psi}}(y)\} \ket{0}.
\]

As a final note, remember that the functional derivatives with respect to the sources, \(\delta/\delta \eta\) and \(\delta/\delta \overline{\eta}\), are also strictly \textbf{anticommuting objects}. This is because differentiation with respect to a Grassmann variable is equivalent to integrating over it. Therefore, when computing correlation functions by taking multiple functional derivatives of \(Z[\eta, \overline{\eta}]\), you must be extremely careful to track the minus signs generated every time a derivative operator passes through another derivative or a fermionic field.

\subsection{A Note on Derivatives of Fields}\TODO{check}

When calculating amplitudes using the path integral formalism, we established that the integral over fields yields the time-ordered vacuum expectation value (the Green's function):
\[
    \langle \Omega | T \{ \hat{\phi}(x_1) \cdots \hat{\phi}(x_n) \} | \Omega \rangle = \int \mathcal{D}\phi \, e^{iS} \phi(x_1) \cdots \phi(x_n).
\]
But what happens if we replace one of the fields in the path integral with an operator containing derivatives, such as $\partial_\mu \phi(x_1)$? This scenario is extremely common when dealing with interaction Hamiltonians containing derivative couplings.

Let us evaluate this explicitly by writing the derivative as a finite difference limit along the time direction, taking $x_1 = (t_1, \vec{x}_1)$:
\[
    \begin{aligned}
        \int \mathcal{D}\phi \, e^{iS} (\partial_{t_1} \phi(x_1)) \phi(x_2) \cdots \phi(x_n) & = \lim_{\epsilon \to 0} \int \mathcal{D}\phi \, e^{iS} \frac{1}{\epsilon} \left[ \phi(t_1+\epsilon, \vec{x}_1) - \phi(t_1, \vec{x}_1) \right] \phi(x_2) \cdots \phi(x_n) \\
                                                                                             & = \lim_{\epsilon \to 0} \frac{1}{\epsilon} \Big[ \langle \Omega | T \{ \hat{\phi}(t_1+\epsilon, \vec{x}_1) \hat{\phi}(x_2) \cdots \} | \Omega \rangle                    \\
                                                                                             & \qquad \qquad \quad - \langle \Omega | T \{ \hat{\phi}(t_1, \vec{x}_1) \hat{\phi}(x_2) \cdots \} | \Omega \rangle \Big]                                                  \\
                                                                                             & = \partial_{t_1} \langle \Omega | T \{ \hat{\phi}(x_1) \hat{\phi}(x_2) \cdots \hat{\phi}(x_n) \} | \Omega \rangle.
    \end{aligned}
\]
Crucially, notice that the derivative operator has factored \textbf{outside} the time-ordered product. In general, it is strictly true that bringing a time derivative inside a time-ordered product changes the result:
\[
    \partial_{t_1} \langle \Omega | T \{ \hat{\phi}(x_1) \cdots \hat{\phi}(x_n) \} | \Omega \rangle \neq \langle \Omega | T \{ \partial_{t_1} \hat{\phi}(x_1) \cdots \hat{\phi}(x_n) \} | \Omega \rangle.
\]

To see why they are not equal, let us calculate a simple example in the canonical operator formalism: the double time derivative of a 2-point function. The time-ordering operator is mathematically implemented using Heaviside step functions $\Theta$:
\[
    \langle 0 | T \{ \phi(x) \phi(y) \} | 0 \rangle = \Theta(x^0 - y^0) \langle 0 | \phi(x) \phi(y) | 0 \rangle + \Theta(y^0 - x^0) \langle 0 | \phi(y) \phi(x) | 0 \rangle.
\]
Applying the first time derivative $\partial_{y^0}$ yields terms from the derivative of the fields, plus terms from the derivative of the step functions (recall that $\partial_{y^0} \Theta(x^0 - y^0) = -\delta(x^0 - y^0)$):
\[
    \begin{aligned}
        \partial_{y^0} \langle 0 | T \{ \phi(x) \phi(y) \} | 0 \rangle = \langle 0 | T \{ \phi(x) \partial_{y^0} \phi(y) \} | 0 \rangle & - \delta(x^0 - y^0) \langle 0 | \phi(x) \phi(y) | 0 \rangle       \\
                                                                                                                                        & + \delta(x^0 - y^0) \langle 0 | \phi(y) \phi(x) | 0 \rangle       \\
        = \langle 0 | T \{ \phi(x) \partial_{y^0} \phi(y) \} | 0 \rangle                                                                & - \delta(x^0 - y^0) \langle 0 | [ \phi(x), \phi(y) ] | 0 \rangle.
    \end{aligned}
\]
Since equal-time fields commute, the commutator $[\phi(\vec{x}), \phi(\vec{y})]$ vanishes, and the delta function term disappears. However, when we apply the second derivative $\partial_{x^0}$, we hit the step functions again:
\[
    \begin{aligned}
        \partial_{x^0} \partial_{y^0} \langle 0 | T \{ \phi(x) \phi(y) \} | 0 \rangle & = \langle 0 | T \{ \partial_{x^0} \phi(x) \partial_{y^0} \phi(y) \} | 0 \rangle + \delta(x^0 - y^0) \langle 0 | [ \phi(x), \partial_{y^0} \phi(y) ] | 0 \rangle.
    \end{aligned}
\]
Now, the equal-time commutator is between the field and its conjugate momentum $\pi(y) = \partial_{y^0} \phi(y)$, which evaluates to $i\delta^{(3)}(\vec{x}-\vec{y})$. Therefore:
\[
    \partial_{x^0} \partial_{y^0} \langle 0 | T \{ \phi(x) \phi(y) \} | 0 \rangle = \langle 0 | T \{ \partial_{x^0} \phi(x) \partial_{y^0} \phi(y) \} | 0 \rangle + i\delta^{(4)}(x - y).
\]
This extra $\delta^{(4)}(x - y)$ is a \textit{contact term} that only exists when the two points perfectly coincide. It mathematically demonstrates why taking derivatives outside the $T$-product is not equivalent to pushing them inside.

\paragraph{Conclusion}
Derivative operations appearing inside the path integral are strictly equivalent to applying the derivative operator to the Green's function, i.e., \textbf{outside the time-ordered product}:
\[
    \int \mathcal{D}\phi \, e^{iS} \left( \frac{\partial}{\partial x_1^\mu} \phi(x_1) \right) \phi(x_2) \cdots = \frac{\partial}{\partial x_1^\mu} \langle \Omega | T \{ \hat{\phi}(x_1) \hat{\phi}(x_2) \cdots \} | \Omega \rangle.
\]
Notice that this perfectly matches what we did when deriving Feynman rules for interaction terms with derivatives. We kept the derivatives outside the Wick contractions (which are T-products) so that they could act directly on the plane waves $e^{\pm i p \cdot x}$, yielding factors of $\pm i p_\mu$. This is the correct, unambiguous prescription when working with the path integral formalism.
\begin{remark}
    In the Hamiltonian formalism, deriving this properly would require adding an extra term that cancels exactly with the difference between $\mathcal{H}_{\text{int}}$ and $-\mathcal{L}_{\text{int}}$, leading to the same physical result.
\end{remark}

\section{Quantization of QED} \TODO{check}

Quantizing gauge theories in the path integral formulation is fundamentally more difficult than quantizing scalar theories because the gauge field $A_\mu(x)$ contains an enormous redundancy. There are infinitely many field configurations that correspond to the exact same physical state (namely, configurations that differ only by a gauge transformation). In the case of QED, the transformation is:
\[
    A_\mu(x) \to A_\mu(x) + \frac{1}{e} \partial_\mu \alpha(x),
\]
for any arbitrary scalar function $\alpha(x)$. Because the action is invariant under this transformation, integrating over all possible configurations $A_\mu$ means we are redundantly integrating over an infinite volume of gauge-equivalent states, causing the path integral to badly diverge.

To see this mathematically, let us look at the free QED action, which can be rewritten (after integration by parts) as:
\[
    S = \int \mathrm{d}^4x \left( -\frac{1}{4} F_{\mu\nu} F^{\mu\nu} \right) = \frac{1}{2} \int \mathrm{d}^4x \, A^\mu \left[ \partial^2 g_{\mu\nu} - \partial_\mu \partial_\nu \right] A^\nu.
\]
In order to compute the propagator (and thus the generating functional), we must invert the differential operator appearing in the quadratic part of the action:
\[
    \hat{O} = -i (\partial^2 g_{\mu\nu} - \partial_\mu \partial_\nu).
\]
The propagator $D_F^{\nu\rho}(x-y)$ would be the Green's function satisfying:
\[
    \hat{O} D_F^{\nu\rho}(x-y) = -i (\partial^2 g_{\mu\nu} - \partial_\mu \partial_\nu) D_F^{\nu\rho}(x-y) = \delta_\mu^{\ \rho} \delta^{(4)}(x-y).
\]
However, this equation has \uline{no solution}. The operator $\hat{O}$ is not invertible because it has a non-trivial null space. Specifically, any pure gauge configuration of the form $A_\mu = \partial_\mu \alpha$ is annihilated by the operator:
\[
    (\partial^2 g_{\mu\nu} - \partial_\mu \partial_\nu) \partial^\nu \alpha = \partial^2 \partial_\mu \alpha - \partial_\mu (\partial^2 \alpha) = 0.
\]
Equivalently, if $A_\mu = \partial_\mu \alpha$, the action is exactly $S=0$, making the path integral wildly divergent.

The exact same issue appears in momentum space. The equation for the momentum-space propagator $\Pi^{\nu\rho}(p)$ is:
\[
    -i (-p^2 g_{\mu\nu} + p_\mu p_\nu) \Pi^{\nu\rho}(p) = \delta_\mu^{\ \rho}.
\]
The matrix $(-p^2 g_{\mu\nu} + p_\mu p_\nu)$ annihilates the vector $p^\nu$, hence it has a determinant of zero and cannot be inverted to find $\Pi^{\nu\rho}$.
To compute a well-defined path integral, we must isolate and remove the integration over these troublesome pure-gauge modes ($A_\mu = \partial_\mu \alpha$) and \uline{integrate only over a subspace of inequivalent, physically distinct field configurations}.

\begin{example}[Massive spin-1 case]
    As an exercise, consider a massive vector field (Proca field). Its Lagrangian includes a mass term $\frac{1}{2}m^2 A_\mu A^\mu$. In momentum space, the operator becomes $(-p^2 g_{\mu\nu} + p_\mu p_\nu + m^2 g_{\mu\nu})$. Because the mass term breaks gauge invariance, this new operator no longer annihilates $p^\nu$. It is invertible, and one can solve for the massive spin-1 propagator: $\Pi^{\mu\nu} = \frac{i}{p^2 - m^2 + i\epsilon} (-g^{\mu\nu} + \frac{p^\mu p^\nu}{m^2})$. This confirms that the non-invertibility in QED is strictly a consequence of exact gauge symmetry.
\end{example}

\subsection{The Faddeev-Popov Method}
To fix the gauge without destroying the structural integrity of the path integral, we employ a brilliant trick developed by Faddeev and Popov. The core idea is to insert a carefully crafted ``identity'' into the path integral that acts as a delta-function, restricting the integration strictly to one representative field configuration per gauge orbit.

First, we choose a gauge fixing condition $G(A) = 0$. For any arbitrary field $A$, we denote its gauge-transformed version as $A_\alpha = A + \frac{1}{e} \partial \alpha$. We want the condition $G(A_\alpha) = 0$ to uniquely define $\alpha$ for any starting field $A$.
We start with the following definition of the identity
\[
    1 = \int \mathrm{d} z_1 \cdots \mathrm{d} z_N \, \delta(g_1(\mathbf{z})) \cdots \delta(g_N(\mathbf{z}))\det\left(\frac{\partial g_i}{\partial z_j}\right),
\]
which can be generalized to a continuum version, namely:
\[
    1 = \int \mathcal{D}\alpha(x) \, \delta(G(A_\alpha)) \det \left( \frac{\delta G(A_\alpha)}{\delta \alpha} \right).
\]
Here, the delta functional enforces the gauge condition at every point in spacetime, and the Faddeev-Popov determinant is the Jacobian of the transformation.

We assume that for our chosen gauge, the determinant is independent of the specific field configuration $A$. This is true for linear gauges like the Lorenz gauge, where $G(A_\alpha) = \partial_\mu A^\mu + \frac{1}{e} \square \alpha$, so the functional derivative $\frac{\delta G(A_\alpha)}{\delta \alpha} = \frac{1}{e} \square$ is just a differential operator independent of $A$.

Inserting this identity into the QED path integral:
\[
    \int \mathcal{D}A \, e^{iS[A]} = \int \mathcal{D}\alpha \det\left( \frac{\delta G}{\delta \alpha} \right) \int \mathcal{D}A \, e^{iS[A]} \delta(G(A_\alpha)).
\]
Since the determinant does not depend on $A$, we can pull it outside. Now, we perform a shift in the integration variable: $A_\mu \to A_\mu - \frac{1}{e}\partial_\mu \alpha$. Because the action is gauge-invariant ($S[A] = S[A_\alpha]$) and the measure $\mathcal{D}A$ is invariant under shifts, the integral over $A$ completely decouples from $\alpha$:
\[
    \int \mathcal{D}A \, e^{iS[A]} = \left( \int \mathcal{D}\alpha \det\left( \frac{\delta G}{\delta \alpha} \right) \right) \int \mathcal{D}A \, e^{iS[A]} \delta(G(A)).
\]
The term in the large parentheses is simply an infinite overall constant (representing the infinite volume of the gauge group), which will cancel out when we normalize by the vacuum partition function.

We now specify a family of gauge conditions: $G(A) = \partial_\mu A^\mu - \omega(x) = 0$, where $\omega(x)$ is an arbitrary scalar function. The determinant remains $\det(\square/e)$. The path integral evaluates to:
\[
    \int \mathcal{D}A \, e^{iS} = (\det (\square/e))\left( \int \mathcal{D}\alpha \right) \int \mathcal{D}A \, e^{iS[A]} \delta(\partial_\mu A^\mu - \omega(x)).
\]
Because this result is mathematically true for \textit{any} choice of $\omega(x)$, it must also be true if we average the right-hand side over all possible functions $\omega(x)$ using a Gaussian weight distribution centered at $\omega = 0$. We multiply by $1 = N(\xi) \int \mathcal{D}\omega \exp\left[ -i \int \mathrm{d}^4x \frac{\omega^2}{2\xi} \right]$ (where $\xi$ is an arbitrary constant parameter) and integrate over $\omega$:
\[
    \int \mathcal{D}A \, e^{iS} = \tilde{N}_0 \int \mathcal{D}A \int \mathcal{D}\omega \, e^{iS[A]} e^{-i \int \mathrm{d}^4x \frac{\omega^2}{2\xi}} \delta(\partial_\mu A^\mu - \omega(x)).
\]
The delta functional trivially allows us to perform the integral over $\omega$, simply replacing every instance of $\omega$ with $\partial_\mu A^\mu$. Absorbing all constant prefactors into $N_0$, we obtain:
\[
    \int \mathcal{D}A \, e^{iS[A]} = N_0 \int \mathcal{D}A \, \exp\left[ i \int \mathrm{d}^4x \left( \mathcal{L} - \frac{1}{2\xi} (\partial_\mu A^\mu)^2 \right) \right].
\]
By using the Faddeev-Popov method, we have effectively added a new ``gauge-fixing'' term to the Lagrangian:
\[
    \mathcal{L}_{\text{G.F.}} = -\frac{1}{2\xi} (\partial_\mu A^\mu)^2.
\]
Crucially, if we want to compute the expectation value of an operator $\hat{O}(A)$ that is strictly \textbf{gauge-invariant}, all the steps above remain valid when replacing $e^{iS} \to e^{iS}O(A)$. The infinite constants $N_0$ perfectly cancel between the numerator and denominator:
\[
    \langle \Omega | T \{ \hat{O}(A) \} | \Omega \rangle = \frac{\int \mathcal{D}A \, e^{i S[A]} O(A)}{\int \mathcal{D}A \, e^{i S[A]}} = \frac{\int \mathcal{D}A \, e^{i \int \mathrm{d}^4x ( \mathcal{L} + \mathcal{L}_{\text{G.F.}} )} O(A)}{\int \mathcal{D}A \, e^{i \int \mathrm{d}^4x ( \mathcal{L} + \mathcal{L}_{\text{G.F.}} )}}.
\]
This procedure smoothly generalizes to full QED (and sQED) where fermions or scalars couple to the photon field, provided the operators we measure are fully gauge-invariant.

\subsection{The Photon Propagator}
With the new gauge-fixing term $\mathcal{L}_{\text{G.F.}}$ added to the action, we can finally invert the quadratic operator and compute the photon propagator.

In momentum space, the addition of $-\frac{1}{2\xi} (\partial_\mu A^\mu)^2$ modifies the operator equation to:
\[
    -i \left[ -p^2 g_{\mu\nu} + \left(1 - \frac{1}{\xi}\right) p_\mu p_\nu \right] \Pi^{\nu\rho}(p) = \delta_\mu^{\ \rho}.
\]
To solve this, we note that the propagator must be a rank-2 Lorentz covariant tensor depending only on the momentum $p$. The most general tensor form it can take is:
\[
    \Pi^{\nu\rho}(p) = A(p^2) g^{\nu\rho} + B(p^2) p^\nu p^\rho.
\]
Inserting this ansatz into the modified operator equation, we can uniquely solve for the coefficients $A$ and $B$. Adding the standard $+i\epsilon$ prescription for the poles, we arrive at the full photon propagator:
\[
    \Pi^{\nu\rho}(p) = \frac{i}{p^2 + i\epsilon} \left( -g^{\nu\rho} + (1 - \xi) \frac{p^\nu p^\rho}{p^2} \right).
\]
Notice that the mathematical form of the propagator depends explicitly on the arbitrary gauge-fixing parameter $\xi$. This is perfectly acceptable because the propagator itself is neither gauge-invariant nor a directly observable physical quantity. When calculating the $S$-matrix elements for actual physical processes, the Ward Identity guarantees that all terms proportional to $p^\nu p^\rho$ will identically vanish, rendering the final physical cross-sections completely independent of $\xi$.

In practice, physicists choose specific values for $\xi$ to maximally simplify calculations:
\begin{itemize}
    \item $\xi = 1$ is the \textbf{Feynman gauge}. This choice brilliantly cancels the second term, reducing the propagator to simply $\frac{-ig^{\nu\rho}}{p^2 + i\epsilon}$. This is the most commonly used gauge in perturbative calculations and is the one we have adopted until now.
    \item $\xi \to 0$ is the \textbf{Lorenz gauge} (often called the \textbf{Landau gauge}). In this limit, the propagator becomes purely transverse ($p_\nu \Pi^{\nu\rho} = 0$). This is a highly popular choice for non-perturbative calculations and lattice QCD.
\end{itemize}